\documentclass[11pt]{article}
% Unicode fonts, because the statement and the prose carry Lean's symbols
% (ℝ, ℚ, ∀, ≠, √, ⟨⟩) and the Latin Modern faces under T1 have no glyph for
% them: Tectonic then logs "Missing character" and emits a PDF whose exact
% Lean statement has silently lost its quantifier and its domains. DejaVu
% covers every symbol the recorded runs used and ships in the pinned bundle.
\usepackage{fontspec}
\setmainfont{DejaVu Sans}
\setmonofont{DejaVu Sans Mono}
\usepackage[margin=1in]{geometry}
\usepackage{fancyvrb}
\usepackage{xcolor}
\usepackage{enumitem}
\usepackage[hidelinks]{hyperref}

\definecolor{HardyBlue}{RGB}{24,64,112}
\setlength{\parindent}{0pt}
\setlength{\parskip}{0.7em}
\setlist[itemize]{leftmargin=1.5em}

\begin{document}

\begin{center}
{\LARGE\color{HardyBlue}\bfseries The Irrationality of √2 + √3\par}
\vspace{1em}
Formal status: Kernel verified\\
Faithfulness: User approved; independently reviewed by claude-opus-5 on claude (agreed)\\
Informal completeness: Not independently assessed\\
Document status: TeX compiled
\end{center}

\hrule
\vspace{1em}

\section*{Theorem}
The real number √2 + √3 is irrational; that is, there is no rational number q with q = √2 + √3.

Formally (Lean 4 / Mathlib):

    theorem irrational\_sqrt\_two\_add\_sqrt\_three :
        Irrational (Real.sqrt (2 : ℝ) + Real.sqrt (3 : ℝ))

Here `Irrational x` is Mathlib's definition `x ∉ Set.range ((↑) : ℚ → ℝ)`, i.e. x is not the image of any rational number under the canonical embedding ℚ ↪ ℝ.

\section*{Approved interpretation}
\begin{itemize}
\item ℝ — the real numbers, Lean/Mathlib's `Real`
\item ℚ — the rational numbers, entering only through Mathlib's `Irrational` predicate, defined as `x ∉ Set.range ((↑) : ℚ → ℝ)`
\item `Real.sqrt : ℝ → ℝ` — the nonnegative square root function on the reals
\item None explicit: the statement is a closed proposition about the single real number `Real.sqrt (2 : ℝ) + Real.sqrt (3 : ℝ)`.
\item One universal quantifier is implicit inside `Irrational`: for every `q : ℚ`, `(q : ℝ) ≠ Real.sqrt 2 + Real.sqrt 3`.
\item No hypotheses; the theorem is unconditional and has no binders.
\item The numerals 2 and 3 are the real numbers 2 and 3, obtained via `OfNat ℝ`; they are ascribed as `(2 : ℝ)` and `(3 : ℝ)` so numeral elaboration is pinned rather than left to unification.
\item `Real.sqrt` is Mathlib's total square-root function, which on nonnegative arguments returns the nonnegative root, so `Real.sqrt 2` and `Real.sqrt 3` are the intended principal roots.
\item `Irrational` and `Real.sqrt` are available from Mathlib (`Mathlib.Analysis.SpecialFunctions.Sqrt` / `Mathlib.RingTheory.Int.Basic` transitively via `import Mathlib`); `Irrational` lives in the root namespace, so it is written unqualified.
\item Repair applied: the binder field is now genuinely empty rather than a parenthetical English note, so the assembled signature is `theorem irrational\_sqrt\_two\_add\_sqrt\_three : Irrational (Real.sqrt (2 : ℝ) + Real.sqrt (3 : ℝ))` with nothing spurious between the name and the colon.
\item Repair applied: explicit `(2 : ℝ)` and `(3 : ℝ)` type ascriptions on the numerals, and fully qualified `Real.sqrt`, to avoid any ambiguity in numeral or name elaboration.
\item 'sqrt(2)' and 'sqrt(3)' are read as the principal (nonnegative) real square roots.
\item 'is irrational' is formalized with Mathlib's `Irrational` predicate rather than an ad hoc `∀ p q : ℤ, q ≠ 0 → ...` statement; the two are equivalent over ℝ.
\item The claim concerns the single sum `Real.sqrt 2 + Real.sqrt 3` (≈ 3.146), not the irrationality of the two summands separately, and asserts nothing about algebraic degree or minimal polynomials.
\item No strengthening or weakening: the logical content is exactly the original claim.
\end{itemize}

\section*{Human-readable proof}
\#\# Overview

The argument is the classical one. A sum of two square roots is hard to attack directly, but squaring it collapses the two roots into a single one: (√2 + √3)² = 5 + 2√6. Rationality is preserved by squaring and by the field operations, so if √2 + √3 were rational then √6 would be rational too. Since 6 is not a perfect square, √6 is irrational, and the contradiction is complete.

\#\# Preliminary: √6 is irrational

We use the standard fact that for a natural number n, the real number √n is rational only when n is a perfect square. For n = 6 this gives irrationality directly, since 2² = 4 < 6 < 9 = 3², so 6 lies strictly between consecutive squares and is therefore not a square.

For completeness, here is the self-contained descent argument for this particular case. Suppose √6 = a/b with integers a, b, b > 0, and with a/b in lowest terms. Then a² = 6b². The right-hand side is even, so a² is even, hence a is even; write a = 2a'. Substituting gives 4a'² = 6b², i.e. 2a'² = 3b². Now the left-hand side is even, so 3b² is even, and since 3 is odd, b² is even, hence b is even. So 2 divides both a and b, contradicting the assumption that a/b was in lowest terms. Therefore √6 is irrational.

In the formal development this step is discharged by Mathlib's `norm\_num` extension for irrationality of square roots (`Mathlib.Tactic.NormNum.Irrational`), which certifies `Irrational (√6)` by exhibiting 6 as strictly between the consecutive squares 2² and 3².

\#\# The main argument

Suppose, for contradiction, that √2 + √3 is *not* irrational. Unfolding the definition, this means there exists a rational number q whose image in ℝ satisfies

    (q : ℝ) = √2 + √3.

We record three elementary facts about the square-root function on the nonnegative reals:

1. (√2)² = 2, since 2 ≥ 0;
2. (√3)² = 3, since 3 ≥ 0;
3. √2 · √3 = √(2 · 3) = √6, using multiplicativity of √ on nonnegative arguments.

Now square the supposed identity. Expanding the binomial and applying (1)–(3):

    (q : ℝ)² = (√2 + √3)²
            = (√2)² + 2·√2·√3 + (√3)²
            = 2 + 2√6 + 3
            = 5 + 2√6.

Rearranging this real-number identity for √6 gives

    √6 = ((q : ℝ)² − 5) / 2.

The crucial point is that the right-hand side is the image of a *rational* number. Indeed, ℚ is a field and the embedding ℚ ↪ ℝ is a ring homomorphism, so the rational number r := (q² − 5)/2 — formed from q by squaring, subtracting 5, and dividing by 2, all inside ℚ — satisfies (r : ℝ) = ((q : ℝ)² − 5)/2 = √6. Hence √6 lies in the range of the embedding ℚ ↪ ℝ, i.e. √6 is rational.

This contradicts the preliminary result that √6 is irrational. Therefore no such rational q exists, and √2 + √3 is irrational. ∎

\#\# Remarks

- The same argument shows more generally that √m + √n is irrational whenever mn is not a perfect square: squaring gives m + n + 2√(mn), so rationality of √m + √n would force rationality of √(mn).
- The hypothesis that *mn* is a non-square is what does the work, not that m and n individually are non-squares. For instance √2 + √8 = 3√2 is still irrational, but that case is not covered by the mn-criterion (mn = 16 is a square) and needs a separate observation.
- The proof never needs the sharper fact that √2 and √3 are individually irrational, nor any theory of field degrees or minimal polynomials; one squaring suffices.

\#\# Formal status

The Lean statement above was rebuilt from the Frozen Claim and checked by a fresh Lean 4 / Mathlib toolchain by Hardy's independent FinalVerifier, and its axiom dependencies were audited: the formal theorem is kernel verified. The formal proof follows the exposition step for step — it introduces the hypothetical rational q by unfolding `Irrational`, supplies (q² − 5)/2 as the rational witness for √6, establishes q² = 5 + 2√6 from `Real.sq\_sqrt` and `Real.sqrt\_mul`, and closes against `Irrational (√6)`.

The English exposition in this document is a human-readable rendering of that formal argument. The prose itself — including the hand-written descent argument for √6 given above, which is an illustrative substitute for the `norm\_num`-certified step rather than a transcription of it — has not been independently machine-checked; only the Lean statement and proof carry that guarantee.

\section*{Known gaps}
\begin{itemize}
\item The English exposition has not itself been machine-checked. Only the Lean theorem and its proof term were rebuilt and kernel verified; the prose is a human-readable rendering and could in principle diverge from the formal text in wording or emphasis.
\item The self-contained parity/descent argument for the irrationality of √6 given in the 'Preliminary' section is presented for the reader's benefit and is not the route taken formally. The formal proof obtains `Irrational (√6)` from Mathlib's `norm\_num` irrationality extension, which certifies it by locating 6 strictly between the consecutive squares 2² and 3². The two arguments establish the same statement, but only the latter is the one that was verified.
\item The generalization mentioned in the Remarks — that √m + √n is irrational whenever mn is not a perfect square — is stated without proof and was not formalized. Only the specific case m = 2, n = 3 is verified.
\item The final rearrangement relies on the standard fact that the coercion ℚ → ℝ is a ring homomorphism, so that (q² − 5)/2 computed in ℚ maps to ((q : ℝ)² − 5)/2 in ℝ. This is used silently in the prose; in Lean it is discharged by `push\_cast` together with linear arithmetic, and is part of the verified proof, but a reader should note it is an appeal to the coercion's algebraic compatibility rather than an evident identification.
\end{itemize}

\section*{Exact Lean statement}
\begin{Verbatim}[fontsize=\small]
theorem irrational_sqrt_two_add_sqrt_three : Irrational (Real.sqrt (2 : ℝ) + Real.sqrt (3 : ℝ))
\end{Verbatim}

\section*{Verification appendix}
\textbf{Axioms used:} propext, Classical.choice, Quot.sound

\begin{Verbatim}[fontsize=\small]
Frozen Claim SHA-256: 8c21bc2ba346c1a4f8dcc7de768d258522a9b51ba91290759377c2b299750dbc
Lean source SHA-256: f329ece01705b287a80e34ef886918d96d5bb1427378fa3a720720c39a9482c2
Verification SHA-256: 328fd2f404128c4c7dc35ddabc5c79c9aba3c1249f7738c23478f2e435dfe261
Prompt set SHA-256: 777f19c97622765dc491943e1d658d959ad114a7301d9d8654a366a36abf5180
Tectonic bundle SHA-256: 6ffe055852f8faf66c0acbe1a7fb27f87b869a90bad1204f3bf4d9683f597c7c
\end{Verbatim}

\begin{samepage}
\textbf{Run and tool identities}
\begin{itemize}
\item Run ID: 7676b533-d3dd-4598-bbf0-477a2cce7b13
\item Model: claude-opus-5
\item Backend: claude
\item Runtime SDK: 0.2.127
\item Lean: 4.33.1
\item Mathlib: 0df444a360eaa60ab8c11dca51a86af692955474
\item Tectonic: 0.16.9
\item Tectonic bundle: https://data1.fullyjustified.net/tlextras-2022.0r0.tar
\end{itemize}
\end{samepage}

\vfill
\small
Lean kernel verification, faithfulness (approved by the user and read back by an
independent model before proof search), heuristic informal review, and document
compilation are independent statuses. This trusted-development MVP
does not sandbox generated Lean or TeX.

\end{document}
